\chapter{Weighting: Rates from the Generated Logic}
\label{ch:weight}

\section{What is being built}

A rewrite rule's left-hand side is a type. Terms inhabit it; the rule fires at its
inhabitants. This chapter is about what happens when one refines that type and attaches
a number to each refinement.

The move is small to state. Given a continued interactive GSLT $G$, the weight
endofunctor $\Wgt_{\Semi}$ produces a theory in which every base rule carries a finite
map
\[
  \varphi_1 \mapsto v_1, \quad \ldots, \quad \varphi_m \mapsto v_m
\]
from formulae refining its left-hand side to values in a semiring $\Semi$. A term's
redexes are then classified by which refinement they satisfy, and the values turn the
set of enabled redexes into a distribution. Where $\Cost$ adjoins a stack of tokens and
$\Hist$ adjoins a word of events, $\Wgt_{\Semi}$ adjoins a table of rates --- and, like
both of them, it adjoins it \emph{to the state}, which is the whole of the technical
content and the source of everything downstream.

Two things distinguish this from the existing literature on stochastic process calculi,
and both are consequences of taking the keys from a logic the theory generates rather
than from a vocabulary the modeller invents.

\paragraph{Keys are formulae, not names.} In Priami's stochastic
$\pi$-calculus~\cite{priami} and in Phillips and Cardelli's Stochastic Pi Machine
(SPiM)~\cite{spim}, a rate is attached to a \emph{channel}. That is one refinement among
many: ``this redex is a communication on $n$'' is a particular formula, and there is no
reason a modeller should be confined to it. Once keys are formulae, a rate may depend on
the shape of the payload, on the \emph{namespace} a channel belongs to --- which in a
reflective calculus is real structure, reachable through the name predicate --- or on any
Boolean combination of these. Channel identity is recovered as the degenerate case, and
Theorem~\ref{wt:thm:spim} says exactly how degenerate it is.

\paragraph{Maps are state, not declaration.} A weight map instance is part of a state
specification. A rule's refinement entries therefore carry not only a value but an
\emph{update function} on the whole map, applied when the rule fires. Rates change during
execution. SPiM does not permit this, and the restriction is not incidental to its design
but constitutive of it, since a static rate function is what makes the underlying object
a chain over terms alone. Markovianity is recovered here by a different route: the map is
\emph{in} the state, so the chain is over configurations $(P, \dec)$ rather than over
terms. That is Theorem~\ref{wt:thm:markov}, and everything in
Chapter~\ref{ch:spiking} depends on it, since plasticity is nothing other than a
weight-map update.

\begin{remark}[Where the keys come from]
\label{wt:rmk:whence}
The formulae used as keys are those of the logic the OSLF construction generates from the
theory's own presentation: one structural connective per term former of $\Sigma$, one
context-labelled modality $\langle K \rangle$ per rewrite rule together with a choice of
redex position, and a propositional layer stratified by the categorical strength the
target specification supplies. That construction is Chapter~\ref{ch:hypercube}, and the
transition system the modalities range over is Chapter~\ref{sec:lts-hml}. Nothing in the
present chapter needs the details; it needs only that the vocabulary of refinements is
\emph{generated with the theory} rather than designed alongside it, so that a modeller who
writes down a language has, at no further cost, the language in which to price it. The
one place the details do intrude is
Observation~\ref{wt:obs:closure}, which is about what a weak target specification cannot
say.
\end{remark}

\section{Keys, and the partition discipline}

Fix a base rule $r : L \rsq R$. Write $L^{\sharp}$ for the structural formula
characterising the left-hand side pattern --- for the rho communication rule, ``an input
on some $n$ in parallel with an output on the same $n$'' with $n$ existentially bound.

\begin{definition}[Refinement]
\label{wt:def:refinement}
A formula $\varphi$ \emph{refines} $L^{\sharp}$ when $\sat \varphi \Rightarrow L^{\sharp}$.
The refinements of $L^{\sharp}$, ordered by entailment, form the \emph{refinement lattice}
of the rule.
\end{definition}

A key set is a finite subset of that lattice. Three requirements are imposed on it, and
the first is the substantial one.

\begin{definition}[Partition discipline]
\label{wt:def:partition}
A key set $\Phi_r = \{\varphi_1, \ldots, \varphi_m\}$ is a \emph{partition of $r$} when
\begin{enumerate}[label=(P\arabic*),leftmargin=3em]
  \item \emph{Exclusivity}: $\sat \neg(\varphi_i \wedge \varphi_j)$ for $i \neq j$;
  \item \emph{Exhaustiveness}: $\sat L^{\sharp} \Rightarrow \bigvee_i \varphi_i$.
\end{enumerate}
\end{definition}

\begin{definition}[Admissible key set]
\label{wt:def:reqs}
A key set is \emph{admissible} when it is a partition and in addition
\begin{enumerate}[label=(R\arabic*),leftmargin=3em]
  \item\label{wt:req:dec} \emph{Decidability}: each $\varphi_i$ lies in a fragment for
        which model checking terminates on the terms of interest;
  \item\label{wt:req:loc} \emph{Locality}: the modal depth of each $\varphi_i$ is bounded
        by some $\alpha_r < \infty$ fixed with the rule;
  \item\label{wt:req:str} \emph{Structurality}: each $\varphi_i$ lies in the structural
        fragment, carrying no behavioural modality $\langle K \rangle$, no greatest fixed
        point, and no hidden-name quantifier.
\end{enumerate}
\end{definition}

\ref{wt:req:dec} is forced by execution: a simulator must classify every enabled redex
before it can sample, so a key whose evaluation may diverge is a key that stops the run.
\ref{wt:req:loc} and \ref{wt:req:str} are forced by \emph{efficient} execution, and they
are restrictions on \emph{keys} alone --- the modalities and the fixed point remain
available for stating and checking properties of the resulting chain. What they may not
do is determine a rate. A rate depending on what a redex can do \emph{next} cannot be
recomputed locally when something elsewhere in the term changes, so an implementation
admitting modal keys must reclassify globally at every step.

\begin{proposition}[Classification]
\label{wt:prop:classify}
Let $\Phi_r$ be a partition of $r$ and let $\Red_r(P)$ be the set of redexes of $r$ in
$P$, each a pair $(k,\sigma)$ of a position and a matching substitution. Then
\[
  \class_r : \Red_r(P) \longrightarrow \Phi_r,
  \qquad
  \class_r(k,\sigma) = \text{the unique } \varphi_i \text{ with } L\sigma \sat \varphi_i
\]
is a total function, for every term $P$.
\end{proposition}

\begin{proof}
$(k,\sigma) \in \Red_r(P)$ gives $P \equiv k[L\sigma]$, so $L\sigma \sat L^{\sharp}$ by
construction of $L^{\sharp}$. Exhaustiveness gives some $i$, which is totality;
exclusivity gives at most one, which is single-valuedness. Well-definedness on
$\equiv$-classes holds because satisfaction is $\equiv$-invariant, the generated
connectives being generated from the signature modulo the equations.
\end{proof}

Proposition~\ref{wt:prop:classify} is the whole reason for the discipline, and what it
buys is visible in what fails without it. Without exhaustiveness the weighting is
\emph{partial}: some enabled redexes have no rate, and one must invent a convention ---
fire them at rate zero, at a default rate, or refuse to run --- that is not part of the
specification. Without exclusivity the weighting is \emph{multi-valued}, and one must
invent an aggregation: take the most specific key if a unique minimum exists, sum the
values of all satisfied keys, take the maximum. Each convention is defensible, and each
yields a different simulator from the same specification, which is precisely the situation
a specification exists to prevent.

\begin{remark}[Completing a partition is cheap]
\label{wt:rmk:default}
The discipline is less onerous than it looks. Given any pairwise-exclusive family
$\{\varphi_1,\ldots,\varphi_m\}$, adjoining
$\varphi_{\bot} := L^{\sharp} \wedge \neg\bigvee_i \varphi_i$ yields a partition, and a
modeller wanting the unlisted redexes inert sets its value to $0$. Exclusivity is the
real constraint; exhaustiveness is a completion. In practice, keys of the form ``the
redex is a communication on channel $n$'', indexed by distinct $n$, are exclusive for
free.
\end{remark}

\subsection{What stating the partition costs}
\label{wt:sec:closure}

Definition~\ref{wt:def:partition} writes (P1) and (P2) as entailments. If they are to be
stated \emph{in the generated logic itself}, that logic must carry negation, conjunction
and disjunction, and a target specification supplying only the finite-limits fragment
will not do --- finite limits supply $\top$ and $\wedge$ and no more
(Remark~\ref{ob:rem:stratification}).

This is a hypothesis and not a theorem about weighting as such. A key family drawn from a
positive fragment may perfectly well \emph{be} pairwise disjoint and jointly exhaustive;
what a weak logic prevents is not the holding of (P1) and (P2) but their \emph{expression}
as entailments between formulae of the language the keys are written in. The alternative
is to discharge them in the metatheory --- by an elaborator checking them on the instances
that arise, or by a proof about a schema of keys --- and the weighting is then perfectly
well defined. Both branches cost something, and the second cost is the interesting one.

\begin{observation}[The cost of an external partition]
\label{wt:obs:closure}
If (P1) and (P2) are discharged in the metatheory rather than in the generated logic, the
specification is no longer \emph{closed}: whether a key set is admissible becomes a
judgment the object language cannot express. Two implementations conforming to the same
written specification may then disagree about admissibility, hence --- by
Proposition~\ref{wt:prop:classify} --- about the classification, hence about the
distribution sampled. That is exactly the failure the partition discipline was introduced
to prevent, relocated one level up rather than removed.
\end{observation}

So a weighted theory must choose between a Boolean target specification and an external
admissibility judgment, and it should say which. This is a small instance of a pattern the
Prestige will make central: a construction that looks self-contained turns out to have a
seam, and the seam is where something the model cannot state has to be supplied from
outside (Chapter~\ref{ch:notation}). It is worth noticing that the seam here is
\emph{visible} --- one can point at the entailment that has to be checked somewhere else
--- and that this is the good case.

There is also a convergence worth recording. Possession of something like a negation and a
conjunction is exactly the condition under which a generated logic enjoys an adequacy
theorem. An internal statement of the partition makes the same demand for an independent
reason: a weighting is a claim about which behaviours are distinguishable enough to be
priced differently, and if that claim is to be made inside the generated logic it needs
the expressive power of the claim that they are distinguishable at all.

\section{The endofunctor}

\begin{definition}[Weight map]
\label{wt:def:wmap}
Let $G$ be a continued interactive GSLT with an admissible key set $\Phi_r$ for each base
rule $r$. A \emph{weight map} is a function
\[
  \dec : \textstyle\coprod_r \Phi_r \longrightarrow \Semi .
\]
Write $\Wmaps$ for the set of weight maps. Each $\Phi_r$ is finite and there are finitely
many rules, so $\Wmaps \cong \Semi^N$ for a fixed $N$.
\end{definition}

\begin{definition}[Configuration]
\label{wt:def:config}
A \emph{configuration} is a pair $\mathfrak{c} = (P, \dec)$ with $P$ a term modulo
$\equiv$ and $\dec \in \Wmaps$. Write $\Cfg$ for the set of configurations.
\end{definition}

The phrase ``a map instance is part of a state specification'' is
Definition~\ref{wt:def:config}, and everything downstream turns on it.

\begin{definition}[Augmented rules]
\label{wt:def:augmented}
An \emph{augmented base rule} is a base rule $r : L \rsq R$ together with, for each
$\varphi_i \in \Phi_r$, an initial value $v_i \in \Semi$ and an \emph{update function}
$u_i : \Wmaps \to \Wmaps$ on the whole map. An \emph{augmented context rule} is a context
rule together with a \emph{geometric factor} $\gfac_c \in \Rnn$ and a fold function
$f_c$ combining the maps arriving from its sibling subterms. The initial weight map is
$\dec_0(\varphi_i) = v_i$.
\end{definition}

\begin{construction}[The weight endofunctor]
\label{wt:con:functor}
$\Wgt_{\Semi}$ sends a continued interactive GSLT $G$, equipped with admissible key sets,
to the theory $\Wgt_{\Semi}(G)$ whose configurations are the pairs of
Definition~\ref{wt:def:config} and whose rules are the augmented rules of
Definition~\ref{wt:def:augmented}: a step is
\[
  (P, \dec) \rewrite (P', \dec'),
  \qquad
  P' \equiv k[R\sigma],
  \qquad
  \dec' = f_{c_1}\bigl(\cdots f_{c_p}(u_i(\dec))\cdots\bigr),
\]
for $(k,\sigma)$ a redex of $r$ with $\class_r(k,\sigma) = \varphi_i$ and $k$ composed of
the context rules $c_1, \ldots, c_p$ from the root inwards, the folds applied from the
redex outwards. In the common case where every $f_c$ projects onto its second argument ---
``the context does not modify the map'' --- this is $\dec' = u_i(\dec)$, and we write it so
below. On morphisms, $\Wgt_{\Semi}$ acts as the identity on the underlying map of theories
and by reindexing on the key sets.
\end{construction}

The parallel with the two preceding chapters is exact, and it is the reason this chapter
sits where it does. Each of the three constructions adjoins a component to the state and
makes the rules act on it: a stack that drains, a word that accrues, a table that is
rewritten. Each leaves the underlying interaction structure alone. And each is a genuine
addition rather than a closure operation --- the table after a step is not in general the
table before it.

\begin{remark}[Two adjunctions and one open question]
\label{wt:rmk:notamonad}
$\Cost$ and $\Hist$ are monads, resolved by free--forgetful adjunctions
(Proposition~\ref{ob:prop:cost-adjunction}, Proposition~\ref{ob:prop:hist-adjunction}).
We claim no such thing for $\Wgt_{\Semi}$. There is an evident candidate unit: the map
sending every key to $1_{\Semi}$, which embeds $G$ as the uniformly weighted theory in
which every enabled redex is equally likely, and which is unital in the right way for the
same reason the empty stack and the empty word are. Since a unit is available it might
well be interesting to ask whether a multiplication exists --- what it would mean to
flatten a weighting of a weighting, and whether the flattening is canonical rather than
chosen. We do not ask it here. Nothing below requires more than an endofunctor, and
stating the construction is the point.
\end{remark}

\begin{remark}[Why context rules carry a factor and not a propensity]
\label{wt:rmk:geometric}
The obvious alternative to Definition~\ref{wt:def:augmented} gives every rule, base and
context alike, a value, and lets the simulator select among all of them. It
double-counts. A step is a base rule firing \emph{somewhere}; the context rules composing
the path from the root to that somewhere are the \emph{address} of the firing, not
additional events. Structural rules are addressing, not dynamics, and incur no primitive
cost. So the context rules contribute multiplicatively, by
$\gfac(k) = \prod_q \gfac_{c_q}$ over the composition, and setting every $\gfac_c = 1$
recovers ``addressing is free'' exactly, which is the default.

The reason to permit $\gfac_c \neq 1$ is that it is the natural home of a \emph{spatial}
bias. $\gfac$ weights an edge of the reduction graph by \emph{where} in the term the
interaction surfaces are; $\dec(\varphi)$ weights it by \emph{what} is being transferred.
That is a geometry/matter split appearing here as an implementation detail, and
Chapter~\ref{sec:discussion} says what would have to be true for it to be more than one.
\end{remark}

\section{Propensity}

The propensity of a rule at a refinement is where the construction is most easily got
wrong, because four independent factors are involved and the literature usually presents
some of them fused.

\begin{definition}[Propensity]
\label{wt:def:propensity}
Let $\mathfrak{c} = (P, \dec)$, let $r$ be a base rule and $\varphi \in \Phi_r$. Put
$K(r,\varphi,P) = \{(k,\sigma) \in \Red_r(P) : \class_r(k,\sigma) = \varphi\}$. The
\emph{propensity} of $(r,\varphi)$ at $\mathfrak{c}$ is
\[
  \prop(r,\varphi,\mathfrak{c})
  \;=\;
  \dec(\varphi) \cdot \!\!\sum_{(k,\sigma) \in K(r,\varphi,P)}\!\! \gfac(k) \cdot \chi(r,k,\sigma),
\]
where $\chi \in \{0,1\}$ is the funding gate of Definition~\ref{wt:def:gate}, identically
$1$ in the absence of cost accounting. The \emph{total propensity} is
$\propz(\mathfrak{c}) = \sum_r \sum_{\varphi \in \Phi_r} \prop(r,\varphi,\mathfrak{c})$.
\end{definition}

The four factors, and what each answers:

\begin{description}[leftmargin=1.6em,style=nextline]
  \item[$\dec(\varphi)$ --- the rate constant.] What kind of interaction this is, priced.
        Dynamic: this is the component the update functions rewrite.
  \item[$\gfac(k)$ --- the geometric factor.] Where the interaction is, priced. Default
        $1$.
  \item[The cardinality of the sum --- multiplicity.] How many ways this interaction is
        available. This is Gillespie's combinatorial count of distinct reactant
        tuples~\cite{gillespie1977exact}, and it is the factor most often omitted in
        implementations, with the consequence that a term with fifty available redexes of
        a kind fires them no faster than a term with one.
  \item[$\chi$ --- the funding gate.] Whether the interaction can be afforded.
\end{description}

Multiplicity deserves a word, because it is the one place where the associative--commutative
structure of the interaction constructor does real work. Redexes are counted as
\emph{selections} from the bag, not identified up to $\equiv$. Two indistinguishable
messages on the same channel offer two ways for a receipt to fire, and a construction that
identified them would run the system at half rate. This is not a subtlety of the
implementation; it is the difference between a truncated Poisson stationary distribution
and a geometric one.

\begin{definition}[Funding gate]
\label{wt:def:gate}
For a redex $(k,\sigma)$ of $r$, let $\Delta(r,k,\sigma)$ be the demand it incurs and
$\Sigma(P,k)$ the supply available at that location, in the sense of
Chapter~\ref{ch:costmonad}. Put $\chi(r,k,\sigma) = 1$ when $\Sigma(P,k)$ funds
$\Delta(r,k,\sigma)$ and $0$ otherwise.
\end{definition}

\begin{proposition}[Cost gates the support, weight distributes over it]
\label{wt:prop:gate}
Unfunded redexes have propensity zero and contribute nothing to $\propz$, while among
funded redexes the distribution is unchanged from the cost-free case up to normalisation.
\end{proposition}

\begin{proof}
$\chi$ appears as a multiplicative factor taking values in $\{0,1\}$; a zero factor
removes the term from the sum, and the remaining terms are untouched.
\end{proof}

Three consequences are worth recording, and they are the reason this chapter belongs
beside Chapter~\ref{ch:costmonad} rather than anywhere else in the book.

\emph{The gate is decidable, so the simulator stays exact.} Were funding a real-valued
penalty rather than a Boolean, the propensity would depend on an optimisation and the
process would no longer be exactly simulable. The decidability of the funding judgment is
doing real work here.

\emph{Exhaustion is a halting condition, not a deadlock.} $\propz = 0$ when every enabled
redex is unfunded. The chain has reached an absorbing state, which is the correct reading:
a system that cannot pay for any of its available transitions has stopped, and the waiting
time to the next event is infinite rather than undefined.

\emph{The two decorations answer different questions.} Cost decides whether a step may
happen; weight decides how fast and how often. Separating them is what lets a model say
that a strong synapse on a starving neuron is silent --- a statement neither decoration
makes alone, and one Chapter~\ref{ch:spiking} needs.

\section{The simulator is a functor out of the image}

$\Wgt_{\Semi}(G)$ is still a rewrite theory: it has terms, equations and rules, and it
steps. Turning it into something one can sample from is a further operation, and the two
should not be run together.

\begin{definition}[The weighted coalgebra]
\label{wt:def:coalgebra}
Let $\Semi^{(-)}$ denote the free-$\Semi$-semimodule construction, $\Semi^{(X)}$ being the
finitely supported functions $X \to \Semi$. Evaluating
Definition~\ref{wt:def:propensity} pointwise gives
\[
  \gamma : \Cfg \longrightarrow \Semi^{(\Cfg)},
  \qquad
  \gamma(\mathfrak{c})(\mathfrak{c}')
  = \!\!\sum_{\substack{(r,k,\sigma) \\ \text{stepping } \mathfrak{c} \to \mathfrak{c}'}}\!\!
    \dec(\class_r(k,\sigma)) \cdot \gfac(k) \cdot \chi(r,k,\sigma).
\]
\end{definition}

The sum runs over \emph{derivations} reaching $\mathfrak{c}'$, not merely over the target.
Assembling the weightings of a fixed theory into a category and taking the Grothendieck
construction gives the fibred form of this assignment, and the simulator is a functor out
of it. We record that to fix the picture and do not need it below.

\begin{remark}[Simulator and interpreter are different objects]
\label{wt:rmk:simulator}
It is worth saying plainly that $\gamma$ is not what a running system does.
An interpreter advances one actual state; a simulator explores how a system \emph{would}
behave under a range of configurations, and it is offline. That difference is what makes
exhaustive construction of the transition graph, model checking, and the complex codomain
of Section~\ref{wt:sec:quantum} viable here and not in an interpreter, and it is why the
endofunctor and the coalgebra are separated above. The endofunctor adjoins the map. The
coalgebra prices it.
\end{remark}

\begin{theorem}[Markovianity on configurations]
\label{wt:thm:markov}
Sampling $\gamma$ in the standard way --- draw a waiting time exponentially with parameter
$\propz(\mathfrak{c})$, select a redex with probability proportional to its contribution,
fire it --- yields a time-homogeneous continuous-time Markov chain on $\Cfg$, with
generator $Q(\mathfrak{c},\mathfrak{c}') = \gamma(\mathfrak{c})(\mathfrak{c}')$ off the
diagonal and $Q(\mathfrak{c},\mathfrak{c}) = -\sum_{\mathfrak{c}' \neq \mathfrak{c}}
Q(\mathfrak{c},\mathfrak{c}')$. It is \emph{not} in general a Markov chain on terms alone.
\end{theorem}

\begin{proof}[Proof sketch]
The propensity is a function of $\mathfrak{c}$ and of nothing else --- in particular not
of the time or the history --- and the successor is determined by $\mathfrak{c}$ together
with the selected redex. That is the definition of a time-homogeneous chain with the
stated generator. For the negative claim it suffices to exhibit two runs reaching the same
term with different maps, which Example~\ref{wt:ex:nonmarkov} does.
\end{proof}

\begin{example}[Dynamic weights are not a term-level phenomenon]
\label{wt:ex:nonmarkov}
Take rho with a single communication rule and keys $\varphi_a, \varphi_b, \varphi_{\bot}$,
where $\varphi_x$ says the redex is a communication on $x$. Let the updates act on
$\dec(\varphi_a)$ alone, by $u_a : w \mapsto w+1$ and $u_b : w \mapsto 2w$, which do not
commute. A term offering exactly one $a$-redex and one $b$-redex, both of which must fire,
reaches the same successor term by either order; starting from $\dec(\varphi_a) = 1$,
firing $a$ first arrives with $(1+1)\cdot 2 = 4$ and firing $b$ first with
$1\cdot 2 + 1 = 3$. The successor carries a single $a$-redex, so the two runs sit at the
same term with expected waiting times $1/4$ and $1/3$. No rate function of the term alone
reproduces both.
\end{example}

Theorem~\ref{wt:thm:markov} is the justification for
Definition~\ref{wt:def:config}, and the shape of the argument is worth being explicit
about, because the design could have gone otherwise. If weight maps could change but were
\emph{not} part of the state, the process on terms would be non-Markovian: the rate out of
$P$ would depend on how $P$ was reached. One would then be in the world of generalised
semi-Markov processes, with no exact simulation algorithm and no chain to model check.
Putting the map in the state restores exactness at the price of a larger state space, and
that price is real --- $\Cfg$ is $\terms \times \Semi^N$, a continuum unless the reachable
weight maps are finite in number.

\begin{remark}[Self-transitions]
\label{wt:rmk:selfloop}
Nothing forbids a rule whose firing returns a configuration $\equiv$-equal to the one it
fired in. Such a redex contributes to $\propz$ but to no off-diagonal entry, so the
diagonal above is $-\sum_{\mathfrak{c}' \neq \mathfrak{c}} Q$ and not $-\propz$. A
self-loop is a fictitious jump in the sense of uniformisation: the two conventions induce
the same distribution of the state at every time and differ only in the number of recorded
events, hence in sojourn statistics. Neither is wrong; mixing them is.
\end{remark}

\subsection{What the restriction to names gives back}

\begin{theorem}[Conservativity over summation-free SPiM]
\label{wt:thm:spim}
Restrict the construction to: the $\pi$-calculus fragment; keys that are channel-identity
predicates and nothing else; static maps, every update the identity; and every geometric
factor $1$. Then the propensity of Definition~\ref{wt:def:propensity} is exactly the
activity computed by the Stochastic Pi Machine for the summation-free
fragment~\cite{spim}, redex multiplicity included. The restriction is strict in each of
the four coordinates.
\end{theorem}

We do not reproduce the proof, which is a computation matching the two definitions of
activity term by term~\cite{StochSim}; the summation correction that the general SPiM
statement would need is empty here, since rho has no summation operator. The theorem is
worth having for what it says about the generality: everything the construction adds
beyond SPiM is visible as the removal of one of four restrictions, and each removal is
independently motivated. Formula-keys remove the first, plasticity the second, spatial
bias the third.

\section{The complex instance: a quantum continuous-time Markov chain}
\label{wt:sec:quantum}

The codomain is a parameter, and the two instances developed here are
$\Semi = \Rnn$, which is everything above, and $\Semi = \CC$. The complex case is not a
relabelling, and this section says what it is. Its physical reading --- and the
obstruction that reading meets --- is Chapter~\ref{sec:quantum}; what follows is the
construction only.

\begin{definition}[Configuration space]
\label{wt:def:hilbert}
Assume the set of reachable weight maps is finite. Let $\Hilb = \ell^2(\Cfg)$ with
orthonormal basis $\{\,|\mathfrak{c}\rangle : \mathfrak{c} \in \Cfg\,\}$ indexed by
\emph{configurations}, not by terms.
\end{definition}

The finiteness assumption is a hypothesis of this section and not a caveat on it:
quantised or saturating weights are mandatory here. And the choice of basis is the same
choice as Definition~\ref{wt:def:config}, made a second time one level up. If the basis
were indexed by terms, a dynamic weight map would change the generator as the system ran,
and there would be no fixed operator to exponentiate.

\begin{definition}[Jump operators]
\label{wt:def:jump}
For a base rule $r$ and $\varphi \in \Phi_r$ with $\dec(\varphi) = z = |z|e^{i\theta}$,
write $\lambda(z) = |z|^2$ and put
\[
  \Lop_{r,\varphi}
  \;=\;
  e^{i\theta}\!\!\!\sum_{\mathfrak{c},\,\mathfrak{c}''}
  \sqrt{\;\sum_{(k,\sigma)\,:\,\mathfrak{c} \to \mathfrak{c}''}
        \lambda(z)\,\gfac(k)\,\chi(r,k,\sigma)\;}\;
  |\mathfrak{c}''\rangle\langle\mathfrak{c}| ,
\]
the inner sum over redexes of $r$ in class $\varphi$ carrying $\mathfrak{c}$ to
$\mathfrak{c}''$.
\end{definition}

Two readings of one formula. Coefficient-wise, the amplitude of a jump channel is the
square root of its aggregate classical rate, times a phase --- which is the only relation
between an amplitude and a rate that dimensional analysis permits, and which is why nothing
need be assumed about $|z|$. Structurally, the targets are \emph{configurations}, so the
update functions appear as edges of this operator rather than as modifications of it.

The square root is outside the sum and not inside, and the choice is forced. Summing
amplitudes over derivations and squaring at the end would make a channel with $m$
indistinguishable derivations carry weight $m^2\lambda$ where the classical construction
carries $m\lambda$; summing rates and taking one square root carries $m\lambda$ exactly.
The selection principle is \emph{conservativity} --- the complex construction must reduce
to the real one on the quantities both compute --- and not physics. Only $\sqrt{\Sigma}$ is
a lift of the classical semantics.

\begin{definition}[Hamiltonian from the equational theory]
\label{wt:def:hamiltonian}
Let $E_h$ be a distinguished subset of the equations of the theory, \emph{not} quotiented
into the basis, presented with real amplitudes $h_e$. Put
\[
  H \;=\; \sum_{e \in E_h} h_e \!\!\sum_{\mathfrak{c} \equiv_e \mathfrak{c}'}\!\!
  \bigl(|\mathfrak{c}'\rangle\langle\mathfrak{c}| + |\mathfrak{c}\rangle\langle\mathfrak{c}'|\bigr),
\]
a Hermitian operator on $\Hilb$. The \emph{quantum continuous-time Markov chain} of the
weighted theory is the pair $(\Hilb, \Lind)$ with
\[
  \Lind(\rho) \;=\; -i[H,\rho]
  \;+\; \sum_{r,\varphi}\Bigl( \Lop_{r,\varphi}\rho\Lop^{\dagger}_{r,\varphi}
  - \tfrac12\bigl\{\Lop^{\dagger}_{r,\varphi}\Lop_{r,\varphi}, \rho\bigr\}\Bigr),
\]
evolving by $\rho(t) = e^{t\Lind}\rho(0)$~\cite{cao2021quantum}.
\end{definition}

Here is the observation this section exists for. A GSLT presents two kinds of relation on
terms --- symmetric, cost-free equations and directed, costed rewrites --- and a quantum
reading has exactly two slots to fill. The rewrites are directed and irreversible, so they
are naturally dissipative and fill the jump operators. The equations are symmetric and
cost-free, so they are naturally unitary and fill the Hamiltonian. But equations
\emph{quotiented into the basis}, as $\equiv$ is above, contribute nothing. Most
presentations quotient everything, so $H = 0$.

\begin{principle}[Where the coherence is]
\label{wt:prin:slogan}
A weighted GSLT is quantum exactly to the extent that its presentation withholds equations
from the quotient.
\end{principle}

That is entirely under the modeller's control, and it says where to look in a presentation
to find out whether complexifying it will do anything at all.

\begin{proposition}[No interference between derivations; coherence only from $H$]
\label{wt:prop:interference}
Under Definition~\ref{wt:def:jump}, distinct redexes of one class with $\equiv$-equal
contracta do not interfere: their rates are summed under the square root, so
$|\langle\mathfrak{c}''|\Lop_{r,\varphi}|\mathfrak{c}\rangle|^2$ is exactly the classical
rate. Distinct classes do not interfere either, their dissipators being summed as quantum
operations. Consequently at $H = 0$ the populations of $(\Hilb,\Lind)$ evolve by the
classical generator of Theorem~\ref{wt:thm:markov}, and every interference effect in the
model is attributable to the commutator $-i[H,\rho]$.
\end{proposition}

\begin{proof}[Proof sketch]
The first claim is Definition~\ref{wt:def:jump} with orthonormality of the basis; the
second is that the dissipator is applied per operator and the results summed. For the
third, read off the diagonal of $\Lind(\rho)$ at $H=0$: the feeding term contributes
$\sum_{\mathfrak{c}'}|\langle\mathfrak{c}|\Lop|\mathfrak{c}'\rangle|^2 p(\mathfrak{c}')$ and
the anticommutator contributes
$-\langle\mathfrak{c}|\Lop^{\dagger}\Lop|\mathfrak{c}\rangle\,p(\mathfrak{c})$, which are the
off-diagonal and diagonal entries of $Q^{\!\top}$ acting on $p$. Off-diagonal elements of
$\rho$ enter neither.
\end{proof}

\begin{theorem}[The Gillespie algorithm is the quantum jump at $H = 0$]
\label{wt:thm:degeneration}
Unravelling $(\Hilb,\Lind)$ into trajectories by the quantum-jump method --- evolve under
the non-Hermitian effective Hamiltonian, jump when the norm crosses a uniform draw ---
reduces at $H = 0$ to the sampling of Theorem~\ref{wt:thm:markov}, exactly. No hypothesis
on the jump structure is needed.
\end{theorem}

So ``Gillespie-inspired'' is a theorem about degeneration and not an analogy. The proof is
a computation with the effective Hamiltonian's norm decay, and we refer to the
note~\cite{StochSim} for it; the hypothesis it once carried was discharged by the
normalisation above, which is a small illustration of how often a correction in one place
pays for itself in another. Two further points from that source are worth recording without
development: a nonzero $H$ does not by itself buy a non-exponential sojourn --- what does is
\emph{variation} in total exit rate across the configurations the coherence connects --- and
the formulae of the generated logic supply, as basis-diagonal projectors, exactly the
atomic propositions that quantum model checking otherwise has to be handed by hand. The
partition discipline is what makes them projectors rather than merely observables.

\section{Where this is used}

The construction is generic over the theory: it applies to anything one can write down as
a language presentation, so one obtains a stochastic machine for the ambient calculus, for
a user-defined domain language, or for cost-accounted rho, by writing the theory rather
than by writing a simulator.

Three places in this book take it up. Chapter~\ref{ch:spiking} runs it at the smallest
scale where the whole apparatus is visible at once, and finds a spiking neural network
with plasticity in it --- synaptic efficacy represented as transition intensity, and the
same rewrite that performs inference performing the learning. Chapter~\ref{sec:weighted}
asks which semiring, and finds that the choice is what separates nondeterminism from cost
from probability from amplitude. And Chapter~\ref{sec:quantum} asks what the complex
instance means, and finds that Principle~\ref{wt:prin:slogan} relocates a problem the
earlier reading could not solve.

% ===================================================================
